\documentclass[11pt]{article}
\usepackage{fontspec}
\usepackage[french]{babel}
\usepackage{amsmath}
\usepackage{amssymb}
\usepackage{amsthm}
\usepackage{mathrsfs}
\usepackage{tikz-cd}
\usepackage{enumitem}

\newcommand{\printedpage}[1]{\marginpar{\footnotesize p.~#1}}
\newcommand{\msp}[1]{\marginpar{\footnotesize [#1]}}
\newcommand{\uHom}{\underline{\mathrm{Hom}}}
\newcommand{\Ind}{\mathrm{Ind}}
\newcommand{\Hom}{\mathrm{Hom}}
\newcommand{\card}{\mathrm{card}}
\newcommand{\Fl}{\mathrm{Fl}}
\newcommand{\Ob}{\mathrm{Ob}}
\newcommand{\Filt}{\mathrm{Filt}}
\newcommand{\hookxrightarrow}[2][]{\lhook\joinrel\xrightarrow[#1]{#2}}

\begin{document}
\setcounter{footnote}{11}

\printedpage{14}
% head: A. Grothendieck : Dérivateurs, Ch. XVIII, Édition des Universités de Montpellier II et Paris VII

d'où
\[
\Fl\,S/x \;\simeq\; \varinjlim_{i} \Fl\,S/x_i \;,
\]
donc
\[
\card\Fl(S/x) \;\leqslant\; (\underbrace{\card I}_{\leqslant\pi})\underbrace{(\sup_{i}\card\Fl(S/x_i))}_{\leqslant\pi} \;\leqslant \pi^2 = \pi \;,
\]

\hfill q.e.d.

\msp{page 76}
% margin: hand-drawn bracket mark in the left margin beside the next two lines
Prouvons c) (\footnote{\textit{Démonstration canulée}, ne marche que si $\pi^{\pi_0} = \pi$\,!!}), i.e. que pour tout $\pi' > \pi \geqslant \pi_0$ et tout $x \in \Ob\Filt^{\pi'}$, on peut écrire
\[
x \;=\; \varinjlim_{I} x_i \qquad I \text{ grande devant } \pi\text{, les } x_i \text{ dans } \Filt^{\pi}\text{, } \card I \leqslant \pi'^{\pi}. \tag*{(4.3.6.5)}
\]
% margin: hand-drawn bracket mark in the left margin beside the next three lines
La démonstration se fait comme dans loc. cit. (p. 151-152), en changeant légèrement la définition de $I$, comme désignant l'ensemble des sous-catégories pleines $i$ de $S/x$ qui sont filtrantes, de cardinal $\leqslant \pi$ \textit{et de plus grandes devant} $\pi_0$ (\footnote{Il n'est pas sûr que $I$ soit grande devant $\pi$, sauf si on suppose $\pi^{\pi_0} = \pi$\,!}). (On pourra supposer $S/x$ filtrante en augmentant $S$ de façon à ce que $S$ soit stable dans $\mathcal{M}$ par petites $\varinjlim$ finies, donc les catégories $S/x$ aussi, ce qui implique qu'elles sont filtrantes. Il est \texttt{[phrase incomplète]})

On a donc prouvé PTA $\Longrightarrow$ c).

\textit{Prouvons c) $\Longrightarrow$ d).} Soit
\[
\pi_1 \;=\; 2^c \qquad \text{avec} \qquad c \;\geqslant\; \pi_0 \;.
\]
On sait alors que
\[
\Filt^{\pi_1}\mathcal{M} \;=\; \mathcal{M}_{\pi_1} \;\overset{\text{déf}}{=}\; \text{sous-catégorie pleine de $\mathcal{M}$ des objets $\pi_1$-accessibles de $\mathcal{M}$}
\]
(loc. cit. 9.16), ce qui implique notamment que $\mathcal{M}_{\pi_1}$ est stable par $\varinjlim$ de cardinal $\leqslant \pi_1$, en vertu de loc. cit. 9.9 (page 145). De plus,
\msp{page 77}
par loc. cit. 9.18 b) (p. 156), on a
\[
\Ind_{\pi_1}(\mathcal{M}_{\pi_1}) \xrightarrow{\;\sim\;} \mathcal{M} \;,
\]
ce qui prouve d).

\textit{Prouvons d) $\Longrightarrow$ PTA.} Par loc. cit. 9.18 a) (p. 156), on voit que les objets de $\mathcal{M}_{\pi_1}$ sont $\pi_1$-accessibles dans $\Ind_{\pi_1}(\mathcal{M}_{\pi_1})$, d'autre part il est clair que $\mathcal{M}_{\pi_1}$ est une sous-catégorie pleine strictement génératrice de $\Ind_{\pi_1}(\mathcal{M}_{\pi_1})$, donc il reste à voir que $\Ind_{\pi_1}(\mathcal{M}_{\pi_1})$ est stable par petites $\varinjlim$ , i.e. le

\textbf{Lemme 4.3.7.} \textit{Soit $\mathcal{M}_\pi$ une $\mathfrak{U}$-catégorie stable par limites inductives de cardinal $\leqslant \pi$, où $\pi$ est un cardinal infini donné. Alors} $\Ind_\pi\mathcal{M}_\pi$ \textit{est stable par petites limites inductives. De plus, $\mathcal{M} = \Ind_\pi\mathcal{M}_\pi$, munie de $\mathcal{M}_\pi \longrightarrow \mathcal{M}$, est solution d'un problème $2$-universel, à savoir
\printedpage{15}
% head: A. Grothendieck : Dérivateurs, Ch. XVIII, Édition des Universités de Montpellier II et Paris VII
que pour toute catégorie $\mathcal{N}$ stable par petites limites inductives, le foncteur restriction induit une équivalence}
\[
\uHom_{!}(\mathcal{M},\mathcal{N}) \xrightarrow{\;\sim\;} \uHom_{!\,\pi}(\mathcal{M}_\pi,\mathcal{N}) \;, \tag*{(4.3.7.1)}
\]
\textit{où le signe ! (resp. $!\,\pi$) désigne la sous-catégorie pleine de la catégorie} $\uHom$\textit{, formée des foncteurs qui commutent aux petites} $\varinjlim$ \textit{(resp. aux limites inductives de cardinal $\leqslant \pi$).}

\msp{page 78}
Prouvons d'abord que les $\varinjlim$ de cardinal $\leqslant \pi$ \textit{dans} $\mathcal{M}_\pi$ sont aussi des $\varinjlim$ \textit{dans} $\mathcal{M} = \Ind_\pi\mathcal{M}_\pi$ (\footnote{Mais \textit{attention}, $\mathcal{M}_\pi \longrightarrow \Ind(\mathcal{M}_\pi)$ ne commute \textit{pas} à ces $\varinjlim$ !}). On a donc
\[
x \xleftarrow{\;\sim\;} \varinjlim_{I} x_i \qquad \textit{dans } \mathcal{M}_\pi \;,
\]
i.e.
\[
\Hom(x,y) \xrightarrow{\;\sim\;} \varprojlim_{I} \Hom(x_i,y) \qquad \text{si } y \in \Ob\mathcal{M}_\pi \;, \tag*{$(*)$}
\]
prouvons que cela reste vrai pour $y$ dans $\mathcal{M}$, i.e. $y$ de la forme
\[
y \;=\; \varinjlim_{J} y_j \;, \qquad J \text{ grande devant } \pi\text{, les } y_j \text{ dans } \mathcal{M}_\pi.
\]
On a en effet
\[
\begin{aligned}
\Hom(x,y) \quad &\simeq \quad \varinjlim_{J} \Hom(x,y_j)\\
&\overset{(*)}{\simeq} \quad \varinjlim_{J}\, \varprojlim_{I} \Hom(x_i,y_j)\\
&\overset{\text{SGA 4 I 9.8}}{\simeq} \quad \varprojlim_{I}\, \varinjlim_{J} \Hom(x_i,y_j)\\
&\simeq \quad \varprojlim_{I} \Hom(x_i,y) \;,
\end{aligned}
\]
au total
\[
\Hom(x,y) \;\simeq\; \varprojlim_{I} \Hom(x_i,y),
\]
q.e.d.

Pour prouver que les petites $\varinjlim$ sont représentables dans $\mathcal{M} = \Ind_\pi(\mathcal{M}_\pi)$, comme elles le sont dans $\Ind(\mathcal{M}_\pi)$ (SGA 4 I 8.9.5 b, on utilise l'hypothèse que les $\varinjlim$ \textit{finies} sont représentables dans $\mathcal{M}_\pi$), il reste à voir que $\Ind_\pi(\mathcal{M}_\pi)$ est stable dans $\mathcal{M}_\pi$ par ces limites.

\msp{page 79}

\section*{4.4\quad Étude de $\Ind_\pi(C)$ (pour $C$ petite catégorie stable par $\varinjlim$ de cardinal $\leqslant \pi$)}

\textbf{4.4.1.} Finalement, la démonstration du lemme 4.3.7 n'est pas tellement évidente que je l'anticipais, et il me faut lui consacrer des préliminaires techniques, que je vais réunir ici. Dans toute la suite, $C$ représente une $\mathfrak{U}$-catégorie équivalente à une petite catégorie.

\printedpage{16}
% head: A. Grothendieck : Dérivateurs, Ch. XVIII, Édition des Universités de Montpellier II et Paris VII

\textbf{Proposition 4.4.2.} \textit{Soit $F \in C^{\wedge}$ un foncteur ind-représentable, i.e. tel que $C/F$ soit filtrante. Soit d'autre part $i \longmapsto x_i$ un foncteur $I \longrightarrow C$, où $I$ est une petite catégorie filtrante, et soit}
\[
\alpha \;:\; x \;=\; \varinjlim_{I}{}_{C^{\wedge}} x_i \longrightarrow F \qquad \text{(\footnotemark)} \tag*{(4.4.2.1)}
\]
\footnotetext{La donnée de $\alpha$ équivaut aussi à celle d'un homomorphisme d'ind-objets $(x_i)_{i\in I} \longrightarrow F$, $F$ étant considéré comme ind-objet (généralisé) indexé par $C/F$.}
\textit{un système projectif de flèches $x_i \longrightarrow F$ (}\textbf{N.B.} \textit{la limite inductive $x$ est prise dans $C^{\wedge}$, et est calculée argument par argument), ce qui équivaut à une factorisation $\varphi_\alpha : I \longrightarrow C/F$ de $I \longrightarrow C$ en}
\[
I \xrightarrow{\;\varphi_\alpha\;} C/F \hookxrightarrow[\psi]{\textit{can.}} C \;.
\]
\textit{Ceci dit, pour que $\alpha$ (4.4.2.1) soit un isomorphisme dans} $\Ind'(C) \subset C^{\wedge}$ \textit{(sous-catégorie pleine des foncteurs ind-représentables sur $C$), il faut et il suffit que $\varphi_\alpha$ soit} cofinal\textit{, ou ce qui revient au même, $I$ étant filtrante, qu'il satisfasse les deux conditions habituelles :}
\begin{itemize}[leftmargin=2.5em]
\item[F\,1)] $\forall\, \xi \in \Ob C/F$, $\exists\, i \in \Ob I$ \textit{et une flèche} $\xi \longrightarrow \varphi_\alpha(i)$.
\item[F\,2)] $\forall\, \xi \in \Ob C/F$, \textit{et} $i \in \Ob I$, \textit{et toute double-flèche}

\msp{page 80}
\[
\xi \rightrightarrows \varphi_\alpha(i) \;,
\]
\textit{il existe une flèche $i \xrightarrow{\;u\;} j$ dans $I$ telle que $\varphi_\alpha(u)$ égalise la double-flèche.}
\end{itemize}
\textsc{Démonstration}. Si $\varphi_\alpha$ est cofinal, on en conclut par définition que $\varinjlim \varphi_\alpha \longrightarrow \varinjlim \psi\varphi_\alpha$ est un isomorphisme, or cette flèche n'est autre que $\alpha$. Inversement, supposons que cette flèche soit un isomorphisme, donc $C/F \xleftarrow{\;\sim\;} C/\varinjlim x_i$, et il reste à voir que le foncteur canonique
\[
I \;\longrightarrow\; C/\varinjlim{}_{C^{\wedge}} x_i
\]
est cofinal, i.e. satisfait F\,1) et F\,2). La vérification est $\pm$ tautologique, en utilisant le fait que pour tout $y$ dans $C$, on a
\[
\Hom(y,\, \varinjlim{}_{C^{\wedge}} x_i) \;\xleftarrow{\;\sim\;}\; \varinjlim{}_{\mathrm{Ens}}\Hom(y,x_i),
\]
et du “calcul” des limites inductives filtrantes dans Ens.

% margin: vertical bar in the left margin beside Corollaire 4.4.3
\textbf{Corollaire 4.4.3.} \textit{Soit $\pi$ un ordinal infini. Pour que $F$ soit dans} $\Ind'_\pi(C)$ \textit{(image essentielle de} $\Ind_\pi(C)$ \textit{dans $C^{\wedge}$), il faut et il suffit qu'il existe une catégorie ordonnée grande devant $\pi$ $I$, et un foncteur cofinal $I \longrightarrow C/F$.}

Il nous faut donc donner des conditions, sur une catégorie $J$ $(= C/F)$, pour qu'il existe une catégorie ordonnée $I$ grande devant $\pi$, et un foncteur cofinal $I \longrightarrow J$.

\msp{page 81}
Une telle catégorie s'appellera \textit{grande devant} $\pi$. (Il est immédiat que dans le cas où $J$ est elle-même ordonnée, cette terminologie est cohérente avec celle déjà adoptée pour les ensembles ordonnés, i.e. que s'il existe une application croissante cofinale $\varphi : I \longrightarrow J$, avec $I$ grand devant $\pi$, alors $J$ est lui-même grand devant $\pi$.) Une telle catégorie est
\printedpage{17}
% head: A. Grothendieck : Dérivateurs, Ch. XVIII, Édition des Universités de Montpellier II et Paris VII
automatiquement filtrante (SGA 4 I 8.1.3 b)). Notons aussi que si on a un foncteur cofinal $J \longrightarrow J'$, si $J$ est grande devant $\pi$, $J'$ l'est aussi (par transitivité de la cofinalité).

\textbf{Proposition 4.4.4.} \textit{Soient $J$ une $\mathfrak{U}$-catégorie, $\pi$ un cardinal infini. Conditions équivalentes :}
\begin{itemize}[leftmargin=4em]
\item[\textit{a)}] \textit{$J$ est grande devant $\pi$, i.e. il existe un ensemble ordonné grand devant $\pi$ et un foncteur cofinal $\varphi : I \longrightarrow J$.}
\item[\textit{b)}] \textit{$J$ satisfait les deux conditions suivantes :}
\begin{itemize}[leftmargin=4em]
\item[$\mathrm{PF}_{\pi}\,1$] $\mathrm{Ord}\,J$ \textit{est grande devant $\pi$, i.e. pour toute famille $(y_\alpha)_{\alpha\in A}$ d'objets de $J$, avec} $\card A \leqslant \pi$\textit{, il existe un majorant $y$ dans $J$ (i.e. un $y$ et des flèches $y_\alpha \xrightarrow{\;v_\alpha\;} y$).}
\item[$\mathrm{PF}_{\pi}\,2$] \textit{Pour $y_0$, $y_1$ dans $J$ et tout ensemble $(u_\alpha)_{\alpha\in A}$ de flèches $y_0 \xrightarrow{\;u_\alpha\;} y_1$ avec} $\card A \leqslant \pi$\textit{, il existe une flèche $y_1 \xrightarrow{\;u'\;} y_2$ qui les égalise, i.e. telle que $u'u_\alpha = u'u_\beta \;\forall\, \alpha,\beta \in A$. (Il suffit de l'exiger quand} $\card A \geqslant 2$.\textit{)}
\end{itemize}
\item[\textit{c)}] \textit{Pour toute partie $A$ de} $\Fl J$\textit{, avec} $\card A \leqslant \pi$\textit{, existe une sous-catégorie $J'$ de $J$, avec} $\card\Fl J' \leqslant \pi$\textit{, telle que $A \subset \Fl J'$ et que $J'$ admette un objet final.}
\item[\textit{d)}] \textit{Soit $I$ l'ensemble des sous-catégories $J'$ de $J$ telles que} $\card J' \leqslant \pi$ \textit{et ayant un objet final}

\msp{page 82}
\textit{et ordonnons $I$ par inclusion. Alors $I$ est grand devant $\pi$, de plus le foncteur canonique $I \longrightarrow J$ $(J' \longmapsto e_{J'} =$ objet final de $J')$ est cofinal} (\footnote{Introduire $I(J)$ avant l'énoncé de la proposition. \textbf{N.B.} La condition de cofinalité $I(J) \longrightarrow J$ n'est pas automatique, comme on voit en prenant $J = B_G$, $G = \{\pm 1\}$.},\footnote{\textbf{N.B.} L'hypothèse de cofinalité est superflue si $J$ est rigide, i.e. tout endomorphisme d'un objet de $J$ est identique.},\footnote{\ N. Éd.\ \ Dans cette condition (d), le “\textit{de plus le foncteur \dots\ est cofinal}” est ajouté par Grothendieck \textit{a posteriori}, ce qui explique sa vaine tentative de l'établir dans l'implication d $\Longrightarrow$ a ci-dessous, implication qui devient évidente après cet ajout. Par ailleurs, les conditions (c) et (d) ne sont \textit{pas} impliquées par les conditions équivalentes (a) et (b). Par exemple, si $J$ est la catégorie ayant un seul objet et un seul morphisme non identique $p$ tel que $p^2 = p$, alors $J$ satisfait aux conditions (a) et (b), mais \textit{pas} aux conditions (c) et (d).}).
\end{itemize}
(\footnote{Autres conditions :
\begin{itemize}[leftmargin=3em]
\item[e)] Pour tout ensemble $A$ avec $\card A \leqslant \pi$, $\mathrm{diag}_J : J \longrightarrow J^A$ est cofinal.
\item[e$'$)] (Si $A_0$ est un ensemble donné, de cardinal $= \pi$.) $J \longrightarrow J^{A_0}$ est cofinal.
\end{itemize}
Voir aussi 4.9.6.9.1 et 4.9.6.9.2, pages 237, 238.})

\textsc{Démonstration}.
\begin{itemize}[leftmargin=7em,labelwidth=6em,align=left]
\item[a $\;\Longrightarrow\;$ b] vérification immédiate, en utilisant le critère F 1, F 2 de la cofinalité pour $I$ filtrante.
\item[b $\;\Longrightarrow\;$ c] par AQT \texttt{[âne qui trotte]}.
\item[c $\;\Longrightarrow\;$ d] par AQT \texttt{[âne qui trotte]}.
\end{itemize}
% foot: Version 17 décembre 2025

\printedpage{18}
% head: A. Grothendieck : Dérivateurs, Ch. XVIII, Édition des Universités de Montpellier II et Paris VII
\begin{itemize}[leftmargin=7em,labelwidth=6em,align=left]
\item[d $\;\Longrightarrow\;$ a] Considérons l'ensemble ordonné grand devant $\pi$ construit dans l'énoncé de d), et le \textit{foncteur}
\[
I \;\longrightarrow\; J
\]
donné par $J' \longmapsto e_{J'}$ $=$ objet final de $J'$. (Vérification que c'est un foncteur immédiate.) Ce foncteur est surjectif, donc satisfait la condition F 1 de cofinalité. La condition F 2 résulte aussi de d) en regardant les trois objets $J'$, $j \xrightarrow{\;u_1\;} e_{J'}$, $j \xrightarrow{\;u_2\;} e_{J'}$ de $I$,
\[
\begin{tikzcd}[column sep=small, row sep=small]
J', & e_{J'} \\
j \arrow[ur, shift left, "u_1"] \arrow[ur, shift right, "u_2"'] &
\end{tikzcd}
\]
et en prenant un $J''$ qui les majore tous trois, alors $e_{J'} \longrightarrow e_{J''}$ égalise $u_1$, $u_2$ (\footnote{Un peu bref, le cas où $e_{J'} = j$ échappe à cet argument, et demande qu'on \texttt{[phrase incomplète]}}).
\end{itemize}

% margin: vertical bar in the left margin beside Corollaire 4.4.5
\textbf{Corollaire 4.4.5.} \textit{Soit $F \in \Ob C^{\wedge}$. Pour que $F$ appartienne à} $\Ind'_\pi(C)$\textit{, il faut et il suffit que $C/F$ satisfasse à l'ensemble des quatre conditions équivalentes de 4.4.4.}

\textbf{Corollaire 4.4.6.} \textit{Soit $J$ une catégorie stable par} $\varinjlim$ \textit{de cardinal $\leqslant \pi$. Alors $J$ est grande devant $\pi$.}

En effet, les conditions $\mathrm{PF}_{\pi}\,1$, $\mathrm{PF}_{\pi}\,2$ sont visiblement vérifiées.

\msp{page 83}
\textbf{4.4.6.} Je me donne un foncteur % [sic] 4.4.6. -> 4.4.7. (?), numéro déjà employé par le Corollaire 4.4.6 qui précède
\[
J \xrightarrow{\;\varphi\;} \Ind'_\pi C \;, \tag*{(4.4.6.1)}
\]
$J$ une petite catégorie, et on cherche une $\varinjlim$ dans $\Ind'_\pi C$, en postulant au besoin l'existence des $\varinjlim$ de cardinal $\leqslant \pi$ dans $C$, et l'hypothèse $\card\Fl J \leqslant \pi$. Le procédé le plus naturel pour ramener la question de l'existence des $\varinjlim$ dans $\Ind'_\pi C$, à la question similaire dans $C$, est explicité dans SGA 4 I 8.8.3. On considère l'inclusion pleinement fidèle
\[
C \;\hookxrightarrow{\;\alpha\;}\; \Ind'_\pi C \;, \tag*{(4.4.6.2)}
\]
laquelle commute aux $\varinjlim$ de cardinal $\leqslant \pi$ (celles qui sont représentables dans $C$, cf. p. 78 et référence à SGA 4 I 9.8), d'où
\[
\uHom(J,C) \;\hookxrightarrow{\;\alpha\;}\; \uHom(J,\Ind'_\pi C) \tag*{(4.4.6.3)}
\]
également pleinement fidèle, et commutant aux $\varinjlim$ de cardinal $\leqslant \pi$ (représentables dans le premier membre). Cette flèche s'insère dans le carré commutatif relatif à la donnée d'un $j \in \Ob J$, et des foncteurs d'inclusion
\[
e \xrightarrow{\;\varepsilon_j\;} J \;, \qquad e_j(e) \;=\; j \;; \tag*{(4.4.6.4)} % [sic] e_j(e) -> \varepsilon_j(e), la flèche est notée \varepsilon_j
\]
% foot: Version 17 décembre 2025

\end{document}
